\section{Methods} \label{sec:method} The method has three parts: an oracle that answers the task
with no model in it, the theory that says when it is exact and what it bounds, and the ladder that uses it.

\subsection{Geometric oracle and identifiability}
\label{sec:oracle}

\textbf{Structures and labels.} A crystal structure is a triple $s = (L, X, Z)$ of lattice, fractional
coordinates $X = \{x_i\}_{i=1}^{N}$ and species. The label map $y(s) = \spg(s)$ is the crystal system returned
by a deterministic symmetry algorithm \citep[spglib,][]{spglib2024} at a fixed tolerance $\tau$, so label
correctness is not a competing explanation for model error.

\textbf{Renders.} The frozen protocol renders a $2\times2\times2$ tiling of the conventional cell through five
orthographic cameras $\mathcal{C} = \{c_v\}_{v=1}^{5}$ (three principal-axis, two oblique) at $768$ px. The
tiling is a legibility choice inherited from the convention. Each camera is a known unit direction $d_v$ with
orthographic projection $\Pi_v(p) = P_v\, p$ for a fixed $2\times3$ matrix $P_v$ with orthonormal rows
perpendicular to $d_v$. The task given to a model is: from $(I_1,\dots,I_5)$, name the crystal system.

\textbf{The oracle.} The oracle answers the same question from the renders' own geometry instead of from
pixels: where classical orthographic factorisation estimates the cameras, the render protocol fixes them, so
we invert them in closed form. It forward-projects the ground-truth atom positions of the conventional cell
through each frozen camera, \emph{discards} the correspondence between views, re-solves it by ray intersection
with cross-view verification, and applies $\spg$ to the reconstruction (Algorithm~S1 of the ESI\dag{}). It is
deterministic and CPU-only, and assumes \emph{ideal extraction}: every atom's centroid is available in every view, including
atoms an image hides, so $R_1$ bounds identifiability from the projection geometry rather than from the
pixels, and a visibility-conditioned variant is a separate measurement we have not built.

\textbf{Two tolerances, and why they must be tied.} The oracle carries a tolerance $\delta$ at
which it accepts a ray intersection and merges the resulting candidates, and the label map carries the
symmetry tolerance $\tau$ at which $\spg$ decides. They interact: at $\delta > \tau$ a merge can combine two
distinct atoms and displace the survivor by an amount $\spg$ can see, so the reconstruction is faithful in
atom count yet wrong in symmetry. \emph{The operating rule is therefore to keep $\delta \le \tau$}, and the
identifiability result below assumes exactly this merge. Throughout $\tau = 0.01$\,\AA{}, and every oracle
number is reported with both values stated. Our released pipeline used $\delta = 0.15$\,\AA{}; we write
$\Rone$ for the ceiling in the exact regime $\delta \le \tau$ and $\RoneRel$ for the same oracle read at that
released tolerance (the governing notation is collected in Section~S1 of the ESI\dag{}).

\begin{definition}[Render ceiling] \label{def:ceiling} For an evaluation sample $\mathcal{S}$ under a frozen
protocol $\mathcal{C}$, the \emph{render ceiling} is $\Rone = |\mathcal{S}|^{-1}\sum_{s \in \mathcal{S}}
\mathbf{1}[\,O(s,\mathcal{C}) = y(s)\,]$. \end{definition}

\begin{definition}[Phantom set]
\label{def:phantom}
Fix $\tau \ge 0$. A species-labelled point $(p,z)$ lies in $\Phi_\tau(s,\mathcal{C})$ if it is further than
$\tau$ from every same-species atom, arises as the intersection of the back-projections of two \emph{distinct}
same-species atom images in some pair of views, and projects within $\tau$ of some same-species atom image in
every view. The set depends on $s$ and $\mathcal{C}$ alone; the componentwise statement accompanies the
proofs.
\end{definition}

\begin{theorem}[Soundness and completeness of the oracle]
\label{thm:ident}
Assume at least two directions in $\mathcal{C}$ are linearly independent and distinct same-species atoms are
more than $2\tau$ apart. Then (a) every atom is recovered exactly, by every independent view pair, so the
accepted set contains the atom set $X$ and does not depend on the enumeration order; (b) every accepted
point further than $\tau$ from all same-species atoms lies in $\Phi_\tau(s,\mathcal{C})$, and conversely
every element of $\Phi_\tau(s,\mathcal{C})$ is accepted; and (c) if $\Phi_\tau(s,\mathcal{C}) = \emptyset$
then merging accepted points at tolerance $\tau$ returns exactly $X$, hence $O(s,\mathcal{C}) = y(s)$.
\end{theorem}

Proofs are in Section~S2 of the ESI\dag{}. Parts (a) and (b) say the oracle has no failure mode other than
$\Phi_\tau$: exact on the atoms, and everything it can get wrong is a coincidence fixed by the structure and
cameras. The separation hypothesis in (c) is not cosmetic, and the statement does not lift to labels, since a
reconstruction containing a phantom may still return $y(s)$.

The hypothesis of (c) is satisfied on this benchmark rather than assumed: a direct census of $\Phi_0$ finds
it empty on every certified structure under the released five-camera protocol, and nonempty once cameras are
removed, so $\Rone = 1.0000$ whenever $\delta \le \tau$ and emptiness is a property of the protocol rather
than of the construction. That census is why we treat identifiability as exact and place the operational
degradation in the extraction tolerance instead.

\begin{proposition}[What the ceiling bounds, and what it does not]
\label{prop:scope}
Let $\kappa = \max_{v \neq w} \|[P_v; P_w]^{+}\|_2$ over the enumerated independent view pairs. Call a
procedure a \emph{centroid-triangulating pipeline} if it extracts species-labelled 2D centroids per view
with error at most $\varepsilon$, triangulates at tolerance $\tau$, and applies $\spg$. At
$\varepsilon = 0$ it returns $O(s,\mathcal{C})$ on every structure, so its accuracy is exactly $R_1$; for
$\varepsilon > 0$ the same holds on every structure none of whose compared distances lies within
$2\kappa\varepsilon$ of $\tau$.
\end{proposition}

The constant is a property of the camera set, $\kappa = 2.7321$ here, and it diverges as two directions
approach parallel, so a protocol of nearly collinear cameras has a ceiling that is fragile to extraction error
even when it is high. Because identifiability here is exact, that fragility is the whole of what distinguishes
protocols on this benchmark, and $\kappa$ is used to
predict it in advance rather than to explain it afterwards. The proposition licenses reading $R_1$ as a
ceiling and fixes the limits: it covers no procedure using shading, occlusion ordering or bond rendering, and
no model that answers without reconstructing. Its band is sufficient for label agreement rather than tight,
which we test empirically.

Two design consequences follow (Corollary~S1 of the ESI\dag{}): adding a view can never destroy identifiability,
and a camera aligned with a lattice vector separates no atoms differing by multiples of it.

Part (b) constrains a single camera and with five it does not bind, so tiling is not an identifiability
question on this protocol but a conditioning one. We prove no rate for how $\Phi_\tau$ grows with atoms per
cell, since the natural union bound needs generic positions and crystal atoms sit on special positions; that
dependence is measured instead and reported as a decomposition.

\subsection{The attribution ladder}
\label{sec:ladder}

\begin{definition}[Rungs and shares]
\label{def:ladder}
For a model $m$ evaluated on the same structures, let
$\Rzero = \tfrac{1}{n}\sum \mathbf{1}[\spg(s) = y(s)]$ (the symmetry algorithm on ground-truth positions),
$\Rone$ the oracle ceiling of Definition~\ref{def:ceiling},
$\Rthree(m)$ the accuracy of $m$ given ground-truth fractional coordinates, species and cell parameters as
text, and
$\Rfour(m)$ the accuracy of $m$ on the pixel renders. Define the \emph{perception component}
$P(m) = \PerceptionComp{m}$, the \emph{post-perception residual} (elsewhere the \emph{symbolic residual})
$S(m) = \PostPerception{m}$, and the \emph{perception share} $P(m) / (R_1 - R_4(m))$. $\Rtwo$ denotes the oracle run
on a learned detector's output.
\end{definition}

$R_1 - R_4(m)$ is a model's total deficit against a model-free reference, $P(m)$ the slice closed by supplying
perception's output as text, and $S(m)$ the slice that supplying it does not close. Because $R_1$ carries no
model in its own measurement, the split is a subtraction against a fixed, certified quantity rather than
against a second model's output. Three conditions make the shares interpretable: the rungs are scored on the
same structures with all comparisons paired; $R_3$ and $R_4$ share a decode budget; and $S(m)$ bounds rather
than isolates symmetry reasoning, a bound we test below and which the test does not clear. Referencing to
$R_1$ rather than $R_0$ costs nothing here, since $\Rzero = \Rone = 1.0000$ and the shares in
Table~\ref{tab:rungs} need no correction; at $\RoneRel$ the two rungs differ by $0.0476$ and every share
would shift downward, which is one practical reason to tie $\delta$ to $\tau$. Note that $R_4$ enters both
numerator and denominator, so the share cannot vary independently of pixel accuracy.

\begin{table*}[t]
\centering
\caption{The attribution ladder ($n = 210$, $K{=}3$). \rlLadderShort{} and \rlWeakestShort{} are
\rlLadder{} and \rlWeakest{}, chosen to span the range. Every derived quantity is a difference of rows above
it. $\Rzero$ is definitional: the label \emph{is} the algorithm applied to the structure. $\Rone$ is
measured, and equal to $\Rzero$: with the extraction tolerance tied to the symmetry tolerance the phantom set
is empty on all 210 structures and the oracle returns every label, so the render protocol withholds nothing
and the entire deficit below is the model's; read at the released merge tolerance this row is
$\RoneRel = 0.9524$. $\Rtwo$, the oracle on a learned detector's output, is unscored under a pre-registered
triangulation-failure gate.}
\label{tab:rungs}
\footnotesize
\setlength{\tabcolsep}{4pt}
\begin{tabular*}{\textwidth}{@{\extracolsep{\fill}}ll S[table-format=1.4] S[table-format=1.4]@{}}
\toprule
& & \multicolumn{2}{c}{Accuracy} \\
\cmidrule(l){3-4}
Rung & Reader & {\small \rlLadderShort} & {\small \rlWeakestShort} \\
\midrule
\multicolumn{4}{@{}l}{\emph{Model-free}} \\
\quad $\Rzero$\; true positions        & symmetry algorithm      & \multicolumn{2}{c}{1.0000} \\
\quad $\Rone$\; render geometry       & inversion, then algorithm & \multicolumn{2}{c}{1.0000} \\[1pt]
\multicolumn{4}{@{}l}{\emph{Model}} \\
\quad $\Rthree$\; exact geometry as text & model                  & 0.8524 & 0.4143 \\
\quad $\Rfour$\; five renders (pixels)  & model                  & 0.7333 & 0.3667 \\
\midrule
\multicolumn{4}{@{}l}{\emph{Derived} (Definition~\ref{def:ladder})} \\
\quad perception component $\Rthree{-}\Rfour$        & & 0.1191 & 0.0476 \\
\quad post-perception residual $\Rone{-}\Rthree$    & & 0.1476 & 0.5857 \\
\quad perception share $P/(P+S)$              & & \bfseries 0.4466 & \bfseries 0.0752 \\
\bottomrule
\end{tabular*}
\end{table*}

\subsection{Benchmark construction}
\label{sec:benchmark}

\textbf{Certified labels.} Each structure's crystal system, Bravais lattice, point group and space group
are derived from its coordinates by $\spg$ at the canonical tolerance ($\Symprec=0.01$, angle
tolerance $5^\circ$), with a tolerance sweep that quarantines any structure whose space group is not stable
across it. On a separate stratified sample disjoint from training and evaluation, our labels agree with the
source-database space group in $220/220$ cases (one-sided exact 95\% lower bound $0.9865$; sampling detail in
Section~S3 of the ESI\dag{}). This is not a check against an independent symmetry-detection method: both source
databases derive their space groups through the same spglib library used here, so this certifies
tolerance-consistency within one algorithm family rather than agreement between distinct algorithms.

\textbf{Source structures.} Structures are drawn from the Materials Project \citep{jain2013commentary} as conventional-cell CIFs under
a composition-exclusion split reserving 13 elements for evaluation only. The released 1,820 structures (1,610 train / 210 eval) have zero
train/eval/prior leakage, under the Materials Project's CC~BY~4.0 licence; the database snapshot is not
recorded.

\textbf{Composition exclusion.} No chemical composition in any evaluation set appears in training. Without
this, symmetry is partly predictable from composition alone, and a model could score well without reading the
image.

\textbf{Label granularity.} The crystal system is the headline label, but the same reconstruction determines
the Bravais lattice, point group and space group, so neither the identifiability result nor the ceiling is
specific to the seven-way target; Section~S4 of the ESI\dag{} scores all four.
